% Introduction — Route B (2026-08-25)
\section{Introduction}
Test-time reinforcement learning for language agents --- updating the
policy from the tasks it encounters during deployment so later tasks are
handled better --- has become a crowded and optimistic space. Systems
report gains on math reasoning~\cite{ttrl2025}, tool
use~\cite{t3rl2026}, instance-specific
refinement~\cite{staror2026}, and live in-episode updates with runtime
serving~\cite{att2026}. The deployment setting is genuinely important:
agents deployed in production receive partial, endogenous feedback and
cannot be retrained offline.

But the setting is also uniquely hard to evaluate honestly. Three
properties conspire:

\begin{itemize}
\item \textbf{Training and serving can silently diverge.} The model
  that receives the gradient may not be the model that generates the
  next first attempt. When they diverge, per-task differences between
  frozen and update arms are unidentifiable: they can come entirely
  from sampling-stream displacement --- update arms draw more random
  numbers --- rather than from policy change.
\item \textbf{The evaluator is entangled with the episode.} Hidden
  evaluators that score first attempts are routinely also used to
  terminate episodes early, select retries, or construct training
  labels. Any of those turns the evaluation signal into a training
  signal.
\item \textbf{Replay and anchors can inject ground truth.} Cross-task
  replay buffers are useful, but if their contents are constructed from
  world internals or hand-labeled correct answers, the ``learned''
  behavior is SFT on leaked labels.
\end{itemize}

We demonstrate, by auditing and progressively repairing one pipeline
(``Agent-TTRL''), that these are not hypothetical concerns. Running the
same experimental protocol at three audit levels produces three
different conclusions:

\begin{enumerate}
\item[F1.] With static serving (the updated model never serves), update
  arms show per-task differences that are pure RNG displacement;
  96/96 task outcomes match across ``different'' update variants.
\item[F2.] After colocating training and serving, but with the hidden
  evaluator controlling early stops and hand-constructed anchors in the
  replay buffer, updates appear to transfer ($+0.070$ AUPC,
  $p{=}0.016$, 7/8 seeds positive). Both effects vanish when the
  leakage is removed.
\item[F3.] Under full protocol isolation (exogenous request RNG,
  observation-only credit, signed replay, atomic shadow commits), the
  same REINFORCE-style updates significantly \emph{harm} performance:
  naive $-0.19$ AUPC ($p{=}0.008$, 0/8 seeds positive), egc $-0.23$
  ($p{=}0.008$) on a 16-task latent-skill stream with Mistral-7B.
  The harm replicates exactly on a second backbone (Qwen2.5-7B:
  naive and egc both 8/8 seeds below a deterministic frozen baseline,
  $p{=}0.008$), and no alternative update rule (verified-success-only,
  DPO-style pair loss, looser gates) transfers either.
\end{enumerate}

The takeaway is not that agent test-time RL cannot work --- it is that
\textbf{positive results in this area must first survive an audit
chain} covering served-policy identity, evaluator isolation, evidence
provenance, and update-commit safety. We contribute that chain as a
reproducible harness, together with a defense that works: a pre-commit
gate validating candidate adapters on accessible-evidence instances
restores AUPC to frozen parity (harmful-commit rate 0/8), and a
frozen baseline that is fully deterministic across seeds under
common-random-numbers pairing. As a positive control, the same serving
runtime deployed against the official \texttt{tau2} benchmark --- with
the benchmark's own agent, user simulator, orchestrator, hidden
evaluator, and LLM judge --- completes both evaluated retail tasks
end-to-end with full reward in 5/10 seeded runs using only a local 14B
model, separating ``the pipeline cannot do the task'' (false) from
``the pipeline cannot safely learn'' (true).
